\documentclass[11pt]{article}
\usepackage[margin=1in]{geometry}
\usepackage{amsmath,amssymb,amsthm,mathtools}
\usepackage{graphicx}
\usepackage[hidelinks]{hyperref}
\usepackage{microtype}

\newtheorem{theorem}{Theorem}
\newcommand{\N}{\mathbb N}
\newcommand{\C}{\mathbb C}
\newcommand{\ellp}{\ell^p}

\title{A strongly forward isometric noncontractive\\
representation of the finite Leavitt algebra}
\author{Automated mathematics run \texttt{fa\_banach\_001}\\
Model: GPT5.6}
\date{9 August 2026}

\begin{document}
\maketitle

\begin{center}
\fbox{\parbox{0.91\linewidth}{\textbf{Status.}
Claimed full negative resolution (counterexample), likely valid pending human
review. The construction works for every finite $d\ge2$ and every
$1\le p<\infty$.}}
\end{center}

\section{Source question}

For $d\ge2$, the Leavitt algebra $L_d$ is the universal unital complex algebra
with generators $s_1,\ldots,s_d,t_1,\ldots,t_d$ satisfying
\begin{equation}\label{eq:leavitt}
 t_js_k=\delta_{jk}1,
 \qquad
 \sum_{j=1}^d s_jt_j=1.
\end{equation}
Phillips calls a representation $\rho:L_d\to L(E)$ \emph{contractive on
generators} if all $\rho(s_j)$ and $\rho(t_j)$ have norm at most one. It is
\emph{strongly forward isometric} if every $\rho(s_j)$ is an isometry and,
for every $\lambda\in\C^d$, the operator
$\rho(s_\lambda)=\sum_j\lambda_j\rho(s_j)$ is a scalar multiple of an
isometry \cite[Definition~2.12]{Phillips}.

After giving a strongly forward isometric representation of $L_\infty$ which
is not contractive on generators, Phillips asks whether the same behavior can
occur for finite $L_d$ \cite[Remark~2.13 and p.~15]{Phillips}.

\begin{figure}[ht]
\centering
\includegraphics[width=.94\linewidth]{figures/open_problem_crop.png}
\caption{The definitions and open question in \cite[p.~10]{Phillips}.}
\end{figure}

\section{Idea of the proof}

The standard representation of $L_d$ on $\ell^p$ identifies $\ell^p$ with
the $p$-direct sum of $d$ copies of itself: its forward generators insert a
vector into disjoint coordinate classes, and its reverse generators extract
those classes. We apply a nonisometric change of basis to the first two copies.

Choose the change of basis so that each new column has $\ell^p$-norm one.
This makes every new forward generator an isometry. More strongly, any linear
combination of the forward generators still has a norm-scaling factor
independent of the input vector, so it is a scalar multiple of an isometry.
But the inverse change of basis has a row of norm greater than one. The
corresponding reverse generator is therefore expansive. The algebraic Leavitt
relations are preserved because the changes on the forward and reverse sides
are inverse matrices.

\section{The counterexample}

\begin{theorem}\label{thm:main}
For every $d\ge2$ and every $1\le p<\infty$, there exists a strongly forward
isometric representation
\[
 \rho:L_d\longrightarrow L(\ell^p(\N))
\]
which is not contractive on generators. Consequently, the finite-$d$ question
in \cite[Remark~2.13]{Phillips} has a negative answer.
\end{theorem}

\begin{proof}
Let $V_j,W_j\in L(\ell^p(\N))$, $1\le j\le d$, be the standard spatial
Leavitt family. Explicitly,
\[
 V_j\delta_n=\delta_{d(n-1)+j},
 \qquad
 (W_j\xi)(n)=\xi(d(n-1)+j).
\]
Thus
\begin{equation}\label{eq:standard-relations}
 W_jV_k=\delta_{jk}I,
 \qquad
 \sum_{j=1}^dV_jW_j=I,
\end{equation}
and the ranges of $V_1,\ldots,V_d$ have disjoint supports. In particular,
for all $a=(a_1,\ldots,a_d)\in\C^d$ and $\xi\in\ell^p$,
\begin{equation}\label{eq:row-norm}
 \left\|\sum_{j=1}^d a_jV_j\xi\right\|_p
 =\left(\sum_{j=1}^d|a_j|^p\right)^{1/p}\|\xi\|_p.
\end{equation}

Set $\alpha=2^{-1/p}$ and define
\begin{align*}
 &S_1=V_1,
 &&S_2=\alpha(V_1+V_2),
 &&S_j=V_j \quad(3\le j\le d),\\
 &T_1=W_1-W_2,
 &&T_2=\alpha^{-1}W_2,
 &&T_j=W_j \quad(3\le j\le d).
\end{align*}
The relations $T_jS_k=\delta_{jk}I$ involving indices at least three follow
immediately from \eqref{eq:standard-relations}. For the first two indices,
\begin{align*}
 T_1S_1&=(W_1-W_2)V_1=I,
 &T_1S_2&=\alpha(W_1-W_2)(V_1+V_2)=0,\\
 T_2S_1&=\alpha^{-1}W_2V_1=0,
 &T_2S_2&=W_2(V_1+V_2)=I.
\end{align*}
Moreover,
\begin{align*}
 S_1T_1+S_2T_2
 &=V_1(W_1-W_2)+(V_1+V_2)W_2\\
 &=V_1W_1+V_2W_2.
\end{align*}
Adding $\sum_{j=3}^dV_jW_j$ gives the identity. Hence the $S_j,T_j$
satisfy \eqref{eq:leavitt}, and the universal property produces a
representation $\rho$ with $\rho(s_j)=S_j$ and $\rho(t_j)=T_j$.

We next verify strong forward isometry. Formula \eqref{eq:row-norm} shows that
$S_1$ is an isometry and, since $2\alpha^p=1$, that $S_2$ is an isometry;
the remaining $S_j$ are already isometries. For
$\lambda=(\lambda_1,\ldots,\lambda_d)\in\C^d$,
\[
 \sum_{j=1}^d\lambda_jS_j
 = (\lambda_1+\alpha\lambda_2)V_1
   +\alpha\lambda_2V_2
   +\sum_{j=3}^d\lambda_jV_j.
\]
Therefore \eqref{eq:row-norm} gives, for every $\xi\in\ell^p$,
\begin{equation}\label{eq:scaling}
 \left\|\sum_{j=1}^d\lambda_jS_j\xi\right\|_p
 =c(\lambda)\|\xi\|_p,
\end{equation}
where
\[
 c(\lambda)=
 \left(
 |\lambda_1+\alpha\lambda_2|^p
 +|\alpha\lambda_2|^p
 +\sum_{j=3}^d|\lambda_j|^p
 \right)^{1/p}.
\]
The triangular change of variables appearing here is invertible, so
$c(\lambda)>0$ whenever $\lambda\ne0$. Dividing the operator in
\eqref{eq:scaling} by $c(\lambda)$ gives an isometry. (For $\lambda=0$ the
operator is the zero scalar multiple of any isometry.) Thus $\rho$ is strongly
forward isometric.

Finally, $W_2$ has norm one, and hence
\[
 \|\rho(t_2)\|=\|T_2\|=\alpha^{-1}=2^{1/p}>1.
\]
Thus $\rho$ is not contractive on generators.
\end{proof}

\section{Matrix interpretation and verification}

On the first two coordinates the construction is the generator change
\[
 B=\begin{pmatrix}1&\alpha\\0&\alpha\end{pmatrix},
 \qquad
 B^{-1}=\begin{pmatrix}1&-1\\0&\alpha^{-1}\end{pmatrix}.
\]
The columns of $B$ have $\ell^p_2$-norm one, which is exactly forward
isometry. The second row of $B^{-1}$ has norm $\alpha^{-1}>1$, which is
exactly the failed reverse-generator contraction. This is also a quick audit
of why both properties can coexist without contradicting the Leavitt
relations.

No numerical or computer-assisted step is used. The proof consists only of
the displayed operator identities and the disjoint-support norm formula
\eqref{eq:row-norm}.

\section{Scope and novelty check}

The construction covers all finite $d\ge2$ and all $1\le p<\infty$, so it
fully answers the existence question as stated. It makes no assertion about
spatiality, closures of representation ranges, or isomorphism classes of the
resulting completed Banach algebras.

A bounded search on 9 August 2026 covered the run registry, solution,
attempt, and proof-gap indexes; exact searches for ``strongly forward
isometric'' with ``contractive on generators''; the verbatim open sentence;
title/citation searches for arXiv:1201.4196; and searches for finite-$d$
Leavitt representations. The search found the source and later general
$L^p$-operator-algebra literature, but no later answer to this exact question.
This supports, but does not certify, novelty.

Recommended human review is high priority but narrow: check the four
first-two-coordinate relations, the reconstruction of the identity, and the
norm formula \eqref{eq:scaling}.

\begin{thebibliography}{1}
\bibitem{Phillips}
N. Christopher Phillips,
\emph{Analogs of Cuntz algebras on $L^p$ spaces},
arXiv:1201.4196v1, 2012.
\end{thebibliography}

\end{document}

